% ==============================================================================
% APPENDIX: FLIGHT CONTROLLER VERSION 1
% ==============================================================================
\section{Appendix: Flight controller version 1}

\begin{frame}{Appendix: flight controller version 1}
    \framesubtitle{MPCC reference commands and PID attitude loops}

    \begin{columns}[T,onlytextwidth]
        \begin{column}{0.53\textwidth}
            \centering
            \resizebox{\linewidth}{!}{%
            \begin{tikzpicture}[
                blk/.style={draw=CSNavy,thick,fill=white,align=center,
                    inner sep=3pt,minimum height=8mm},
                pid/.style={draw=red!70!black,thick,fill=white,
                    minimum width=12mm,minimum height=7mm},
                sum/.style={circle,draw=red!70!black,thick,fill=white,
                    minimum size=5mm,inner sep=0pt},
                sig/.style={->,thick,CSNavy},
                fb/.style={->,thick,red!70!black},
                >={Stealth[length=2mm]},
                hop/.style={preaction={draw=white,line width=3pt}}]
                \node[blk,minimum width=17mm,minimum height=26mm] (mpcc)
                    at (0,0) {\scriptsize MPCC\\[-2pt]\scriptsize RNN model};
                \node[sum] (s1) at (2.2,1.0) {};
                \node[sum] (s2) at (2.2,0) {};
                \node[sum] (s3) at (2.2,-1.0) {};
                \node[pid] (p1) at (4.1,1.0) {\scriptsize PID$_\phi$};
                \node[pid] (p2) at (4.1,0) {\scriptsize PID$_\theta$};
                \node[pid] (p3) at (4.1,-1.0) {\scriptsize PID$_r$};
                \node[blk,minimum width=10mm,minimum height=26mm] (mix)
                    at (6.0,0) {$\bar M^{-1}$};
                \node[blk,minimum width=20mm,minimum height=26mm] (drone)
                    at (8.5,0) {\scriptsize Drone\\[-2pt]\scriptsize Simulink};

                \draw[sig] (mpcc.east |- s1) -- node[above,font=\scriptsize]
                    {$\phi_{\rm ref}$} (s1);
                \draw[sig] (mpcc.east |- s2) -- node[above,font=\scriptsize]
                    {$\theta_{\rm ref}$} (s2);
                \draw[sig] (mpcc.east |- s3) -- node[above,font=\scriptsize]
                    {$r_{\rm ref}$} (s3);
                \draw[sig] (s1) -- (p1);
                \draw[sig] (s2) -- (p2);
                \draw[sig] (s3) -- (p3);
                \draw[sig] (p1) -- (mix.west |- p1);
                \draw[sig] (p2) -- (mix.west |- p2);
                \draw[sig] (p3) -- (mix.west |- p3);
                \draw[sig] (mix) -- (drone);
                \draw[fb,hop] (drone.east |- s1) -- (9.8,1.0) -- (9.8,-2.0)
                    -- (2.2,-2.0) -- (s1.south);
                \draw[fb,hop] (9.8,1.0) -- (s2.south);
                \draw[fb,hop] (9.8,1.0) -- (s3.south);
                \node[font=\scriptsize,red!70!black] at (6.2,-2.0)
                    {measured attitude angles};
            \end{tikzpicture}}
        \end{column}

        \begin{column}{0.45\textwidth}
            \begin{block}{Architecture}
                The MPCC computes references for roll, pitch, yaw rate and
                thrust. Each attitude reference was tracked directly by a
                PID controller, before the motor mixing stage.
            \end{block}

            \begin{block}{Limitation}
                This first architecture was stable, but the angle-only
                feedback did not provide enough responsiveness. It was
                therefore replaced by the cascade controller presented in
                the main section, with a fast angular-rate inner loop.
            \end{block}
        \end{column}
    \end{columns}
\end{frame}
